\documentclass{report}

\usepackage[utf8]{inputenc}
\usepackage[spanish]{babel}
\usepackage{geometry}
\usepackage{graphicx}
\usepackage{titlesec}
\usepackage{lipsum}
\usepackage[dvipsnames]{xcolor}
\usepackage[fleqn]{mathtools}
\usepackage{booktabs}
\usepackage{amsmath}
\usepackage{amssymb} % Para los símbolos de conjuntos de números
\usepackage{latexsym}
\usepackage{nccmath}
\usepackage{multicol}
\usepackage{listings}
\usepackage{tasks}
\usepackage{color}
\usepackage{float}
\usepackage{enumitem}
\usepackage{longtable}
\usepackage{makecell}
\usepackage{caption}
\usepackage[parfill]{parskip}
\usepackage{lipsum}
\usepackage{enumitem}
\usepackage[hidelinks]{hyperref}% Para manejar enlaces
\usepackage{fancyhdr}
\usepackage{chappg}
\usepackage{adjustbox}
\usepackage{tocloft}
\usepackage{soul}
\usepackage{pgffor}
\usepackage{tikz}
\usetikzlibrary{arrows}
\usepackage{pgfplots}
\pgfplotsset{compat=1.17}
%Definicion de colores
\definecolor{colorIPN}{rgb}{0.5, 0.0,0.13}
\definecolor{colorESCOM}{rgb}{0.0, 0.5,1.0}
\definecolor{myorange}{RGB}{255, 204, 153}
\graphicspath{ {imagenes/} }
%Configuraciones extras de colores
\definecolor{codegreen}{rgb}{0,0.6,0}
\definecolor{codered}{rgb}{0.6,0,0}
\definecolor{codegray}{rgb}{0.5,0.5,0.5}
\definecolor{codepurple}{rgb}{0.58,0,0.82}
\definecolor{backcolour}{rgb}{0.95,0.95,0.92}

% Configuración de la portada
\titleformat{\section}[block]{\normalfont\huge\bfseries}{\thesection}{1em}{}
\geometry{a4paper, margin=1in}

% Definir comandos para información repetitiva
\newcommand{\logoInstitucion}{recursos/logotipo_ipn.png} % Reemplaza con el nombre del archivo de tu logo
\newcommand{\logoUniversidad}{recursos/EscudoESCOM.png} % Reemplaza con el nombre del archivo de tu logo
\newcommand{\nombreInstituto}{Instituto Politecnico Nacional}
\newcommand{\facultad}{Escuela Superior de Computo}
\newcommand{\materia}{Tópicos Selectos de Algoritmos Bioinspirados}
\newcommand{\grupo}{7BM1}
\newcommand{\profesora}{Daniel Molina Pérez}
\newcommand{\periodo}{2025/02}

\newcommand{\alumnoA}{Carrillo Barreiro José Emiliano}
\newcommand{\alumnoB}{Martinez Ayala Gerardo}
\newcommand{\alumnoC}{Robles Otero José Ángel}
\newcommand{\alumnoD}{Vásquez Morales Haniel Ulises}
\newcommand{\alumnos}{Carrillo J., Martinex G., Robles J. \& Vásquez H.}

% Definir nuevas listas para información repetitiva
\newlist{bracketed}{enumerate}{1}
\setlist[bracketed,1]{label={[{\arabic*}]},left=0pt}

% Configuración del encabezado y pie de página
\pagestyle{fancy}
\fancyhf{} % Limpia todos los campos

% Encabezado
\fancyhead[C]{\textbf{\nombreInstituto\\ \facultad}}
\fancyhead[L]{\includegraphics[width=1cm]{recursos/logotipo_ipn.png}}
\fancyhead[R]{\includegraphics[width=1.5cm]{recursos/EscudoESCOM.png}}

% Pie de página
\fancyfoot[C]{\materia||\grupo\\ \alumnos\\ \thepage}
% Configuración del grosor de la línea en el encabezado y pie de página
\renewcommand{\headrulewidth}{0.4pt}
\renewcommand{\footrulewidth}{0.4pt}

% Ajuste del tamaño del encabezado
\setlength{\headheight}{34.0845pt}
\addtolength{\topmargin}{-22.0845pt}

\fancypagestyle{plain}{%
    \fancyhf{} % Borra encabezado y pie de página por defecto
    \fancyhead[C]{\textbf{\nombreInstituto\\ \facultad}}
    \fancyhead[L]{\includegraphics[width=1cm]{recursos/logotipo_ipn.png}}
    \fancyhead[R]{\includegraphics[width=1.5cm]{recursos/EscudoESCOM.png}}
    \fancyfoot[C]{\materia||\grupo\\ \alumnos\\ \thepage}
    \renewcommand{\headrulewidth}{0.4pt}
    \renewcommand{\footrulewidth}{0.4pt}
}


% Configuración de lstlisting para Python
\lstdefinestyle{mystyle}{
    backgroundcolor=\color{backcolour},
    commentstyle=\color{codegreen},
    keywordstyle=\color{codepurple},
    numberstyle=\tiny\color{codegray},
    stringstyle=\color{codered},
    basicstyle=\ttfamily\small,
    breakatwhitespace=false,
    breaklines=true,
    captionpos=b,
    keepspaces=true,
    numbers=left,
    numbersep=5pt,
    showspaces=false,
    showstringspaces=false,
    showtabs=false,
    tabsize=2
}

\lstdefinestyle{custompython}{
    language=Python,
    basicstyle=\small\ttfamily,
    keywordstyle=\color{blue},
    stringstyle=\color{green},
    commentstyle=\color{red},
    numbers=left,
    numberstyle=\tiny\color{gray},
    stepnumber=1,
    numbersep=5pt,
    backgroundcolor=\color{white},
    frame=single,
    captionpos=b,
    breaklines=true,
    breakatwhitespace=false,
    showstringspaces=false,
    morekeywords={import,plt,np,ws,y,gradient}
}

% Definición de estilo para C++
\lstdefinestyle{cppstyle}{
    language=C++,
    basicstyle=\small\ttfamily,
    keywordstyle=\color{blue},
    commentstyle=\color{green!40!black},
    stringstyle=\color{orange},
    showstringspaces=false,
    breaklines=true,
    breakatwhitespace=true,
    numbers=left,
    numberstyle=\tiny,
    stepnumber=1,
    numbersep=5pt,
    frame=single,
    captionpos=b,
    backgroundcolor=\color{gray!5},
    xleftmargin=0.5cm,
    xrightmargin=0.5cm
}

% Definición de estilo para MATLAB
\lstdefinestyle{matlabstyle}{
    language=Matlab,
    basicstyle=\small\ttfamily,
    keywordstyle=\color{blue},
    commentstyle=\color{green!40!black},
    stringstyle=\color{purple!50!blue!50!white},
    showstringspaces=false,
    breaklines=true,
    breakatwhitespace=true,
    numbers=left,
    numberstyle=\tiny,
    stepnumber=1,
    numbersep=5pt,
    frame=single,
    captionpos=b,
    backgroundcolor=\color{gray!5},
    xleftmargin=0.5cm,
    xrightmargin=0.5cm
}

%Definicion de estilo para Python
\lstdefinestyle{pythonstyle}{
    language=Python,
    basicstyle=\small\ttfamily,
    keywordstyle=\color{red!80!black},
    commentstyle=\color{gray!40!white},
    stringstyle=\color{yellow!30!orange!80!white},
    showstringspaces=false,
    breaklines=true,
    breakatwhitespace=true,
    numbers=left,
    numberstyle=\tiny,
    stepnumber=1,
    numbersep=5pt,
    frame=single,
    captionpos=b,
    backgroundcolor=\color{gray!5},
    xleftmargin=0.5cm,
    xrightmargin=0.5cm
}

% Personalización del índice de listados
\newcommand{\listcppname}{Índice de Listados de C++}
\newcommand{\listmatlabname}{Índice de Listados de MATLAB}
\newcommand{\listpythonname}{Índice de Listados de Python}
\newlistof{cpplistings}{lolcpp}{\listcppname}
\newlistof{matlablistings}{lolmatlab}{\listmatlabname}
\newlistof{pythonlistings}{lolpython}{\listpythonname}

% Configuración para los listados de C++
\lstset{style=cppstyle}
\newcommand{\cpplisting}[1]{%
    \refstepcounter{cpplistings}%
    \addcontentsline{lolcpp}{cpplistings}{\protect\numberline{\thecpplistings}#1}\par}

% Configuración para los listados de MATLAB
\lstset{style=matlabstyle}
\newcommand{\matlablisting}[1]{%
    \refstepcounter{matlablistings}%
    \addcontentsline{lolmatlab}{matlablistings}{\protect\numberline{\thematlablistings}#1}\par}

% Configuración para los listados de Python
\lstset{style=pythonstyle}
\newcommand{\pythonlisting}[1]{%
    \refstepcounter{pythonlistings}%
    \addcontentsline{lolpython}{pythonlistings}{\protect\numberline{\thepythonlistings}#1}\par}

\begin{document}

% Portada
% portada
\begin{titlepage}
    \begin{center}
        \vspace*{1cm}

        \begin{tabular}{c@{\hspace{2cm}}c}
            \includegraphics[width=0.25\textwidth]{\logoInstitucion} &
            \includegraphics[width=0.4\textwidth]{\logoUniversidad}
        \end{tabular}

        \vspace{1cm}

        \textbf{\LARGE \nombreInstituto} \\
        \textbf{\Large \facultad} \\
        \vspace{0.3cm}
        \textbf{\large Materia: \materia} \\
        \textbf{\large Grupo: \grupo} \\
        \vspace{0.3cm}
        \textbf{\large Profesor: \profesora} \\
        \textbf{\large Periodo: \periodo} \\

        \vspace{0.5cm}

        \textbf{\LARGE Practica 02} \\
        \vspace{0.5cm}
        \textbf{\Large \textit{\textbf{Optimización de Carteras:} Un Enfoque de Algoritmo Genético con Manejo de Restricciones para la Selección de Activos}} \\

        \vspace{0.3cm}

        \textbf{\large Realizado por:} \\
        \textbf{\large \alumnoA\\ \alumnoB\\ \alumnoC\\ \alumnoD}

        % Bloque de Resumen y Abstract
        \begin{minipage}{0.8\textwidth}
            \textbf{Abstract:}\\[0.15cm]
            This paper presents an approach to solve the portfolio optimization problem using a genetic algorithm with constraint handling. The goal of portfolio optimization is to find the optimal allocation of capital among different financial assets to maximize return and/or minimize risk. The genetic algorithm is implemented to address this problem, incorporating constraint handling through the exterior penalty method to ensure that solutions satisfy real-world limitations, such as investment limits per asset and acceptable risk levels. The results obtained from applying the algorithm to a stock selection case study are presented and discussed.
        \end{minipage}

        \begin{minipage}{0.8\textwidth}
            \textbf{Resumen:}\\[0.15cm]
            Esta práctica aborda el problema de la optimización de carteras de inversión, donde el objetivo es determinar la distribución óptima del capital entre diferentes activos financieros para maximizar el retorno y/o minimizar el riesgo. Se implementa un algoritmo genético para resolver este problema, incorporando el manejo de restricciones mediante el método de penalización exterior para asegurar que las soluciones obtenidas cumplan con las limitaciones del mundo real, como límites de inversión por activo y niveles de riesgo aceptables. Se presentan y discuten los resultados obtenidos al aplicar el algoritmo a un caso de estudio de selección de acciones.
        \end{minipage}

        \vspace{0.3cm}

        \textbf{\large Fecha: \today}

    \end{center}
\end{titlepage}


% Índice
\tableofcontents
\newpage
%\listofcpplistings%Crea indice de codigos en C++
%\listofmatlablistings%Crea indice de codigos en MATLAB
\listoftables
\listoffigures
%\listoftables

%Capitulos
\chapter{Introducción}
Los algoritmos genéticos (AG) son una clase de algoritmos bioinspirados ampliamente utilizados para resolver problemas de optimización y búsqueda. Inspirados en el proceso de evolución natural, estos algoritmos imitan mecanismos biológicos como la selección, el cruzamiento y la mutación para explorar el espacio de soluciones y encontrar resultados óptimos. En este proyecto se implementa un AG para la minimización de funciones benchmark, en particular las funciones Langermann y Drop-Wave, con el objetivo de demostrar y analizar la efectividad de técnicas como la selección por torneo, el cruzamiento Simulated Binary Crossover (SBX) y la mutación polinomial.

\chapter{Objetivos del Proyecto}

\section{Objetivo general}
Desarrollar e implementar un algoritmo genético para la resolución de un problema de optimización restringida en el ámbito financiero, con el fin de determinar la cartera de inversión óptima bajo diferentes criterios de maximización de retorno y control de riesgo.

\section{Objetivo del inciso uno}
Determinar la asignación óptima de capital en una cartera compuesta por seis acciones, maximizando el retorno esperado bajo la restricción de no invertir más del 40% en una sola acción, sin considerar el riesgo asociado.

\section{Objetivo inciso dos}
Diseñar una estrategia de inversión que minimice el riesgo asociado a la cartera, garantizando un rendimiento esperado mínimo del 35% y respetando la condición de no invertir más del 40% en una sola acción.

\section{Objetivo inciso tres}
Identificar la combinación de inversión que maximice el rendimiento esperado de la cartera sin sobrepasar un nivel de riesgo del 20%, manteniendo la restricción de concentración del 40% por acción.

\chapter{Metodolog\'ia}

\section{Inicializaci\'on de la Poblaci\'on}
La poblaci\'on inicial se genera de forma uniforme a lo largo del espacio de b\'usqueda, definido por l\'imites inferiores y superiores para cada variable.

\subsection*{Objetivo}
Garantizar que la b\'usqueda comience explorando de manera equitativa todas las regiones posibles, evitando sesgos que puedan limitar la diversidad de soluciones iniciales.

\subsection*{Implementaci\'on}
Se utiliza la funci\'on \texttt{initialize\_population}, la cual emplea m\'etodos de generaci\'on aleatoria (por ejemplo, la funci\'on \texttt{np.random.uniform} de NumPy) para crear un conjunto de individuos.

\subsection*{Ventajas}
\begin{itemize}
    \item Permite cubrir todo el rango definido para cada variable.
    \item Aumenta la probabilidad de encontrar regiones prometedoras del espacio de soluciones desde el inicio.
\end{itemize}

\section{Evaluaci\'on de Fitness}

Cada individuo generado se eval\'ua mediante la funci\'on objetivo, la cual determina qu\'e tan buena es la soluci\'on propuesta.

\subsection*{Funciones Utilizadas}
\begin{itemize}
    \item \textbf{Langermann}: Es una funci\'on multimodal que combina componentes cosenoidales y exponenciales, generando m\'ultiples \'optimos locales.
    \item \textbf{Drop-Wave}: Una funci\'on bidimensional con una superficie ondulada, usada para analizar el comportamiento del algoritmo en entornos con m\'ultiples picos y valles.
\end{itemize}

\subsection*{Proceso}
Se calcula el valor de fitness para cada individuo (por ejemplo, evaluando $f(x_1, x_2)$) y se almacena dicho valor para posteriores comparaciones.

\subsection*{Importancia}
La evaluaci\'on correcta del fitness es crucial, ya que determina la selecci\'on de individuos y, por ende, el rumbo de la evoluci\'on poblacional.

\section{Selecci\'on por Torneo}

Para elegir los padres que generar\'an la siguiente generaci\'on se utiliza un m\'etodo de selecci\'on por torneo.

\subsection*{Mecanismo}
\begin{itemize}
    \item Se forman m\'ultiples grupos (torneos) de individuos seleccionados al azar.
    \item En cada grupo se compara el fitness de los participantes y se selecciona al individuo con el mejor desempe\~no.
\end{itemize}

\subsection*{Implementaci\'on Vectorizada}
La funci\'on \texttt{vectorized\_tournament\_selection} realiza este proceso de forma eficiente, aprovechando operaciones vectorizadas de NumPy.

\subsection*{Beneficios}
\begin{itemize}
    \item Favorece la selecci\'on de soluciones de alta calidad sin descartar por completo la diversidad poblacional.
    \item Permite controlar la presi\'on selectiva mediante el tama\~no del torneo.
\end{itemize}

\section{Cruzamiento con SBX}

El operador de cruzamiento se implementa mediante el m\'etodo SBX (Simulated Binary Crossover).

\subsection*{Proceso del SBX}
\begin{itemize}
    \item A partir de dos padres, se genera un n\'umero aleatorio $u$ y se calcula un par\'ametro $\beta$ que determina la dispersi\'on de los descendientes respecto a los padres.
    \item Se generan dos hijos combinando linealmente los valores de los padres.
\end{itemize}

\subsection*{Ajuste de L\'imites}
Se incorpora un mecanismo en \texttt{sbx\_crossover\_with\_boundaries} que garantiza que los hijos resultantes se mantengan dentro de los l\'imites predefinidos.

\subsection*{Ventajas}
\begin{itemize}
    \item Promueve la creaci\'on de soluciones intermedias que pueden explotar la informaci\'on gen\'etica de ambos padres.
    \item Ayuda a preservar la diversidad en la poblaci\'on.
\end{itemize}

\section{Mutaci\'on Polinomial}

Para introducir variabilidad y explorar nuevas regiones del espacio de b\'usqueda, se aplica la mutaci\'on polinomial.

\subsection*{Mecanismo de la Mutaci\'on}
\begin{itemize}
    \item Cada gen de un individuo tiene una probabilidad definida de sufrir una mutaci\'on.
    \item Se usa una distribuci\'on polinomial, controlada por el par\'ametro $\eta_{\text{mut}}$.
\end{itemize}

\subsection*{Consideraciones de L\'imites}
La mutaci\'on se aplica respetando los l\'imites definidos para cada variable mediante la funci\'on \texttt{polynomial\_mutation\_with\_boundaries}.

\subsection*{Beneficios}
\begin{itemize}
    \item Introduce peque\~nas variaciones que pueden conducir a la exploraci\'on de nuevas soluciones.
    \item Previene la convergencia prematura al mantener la diversidad gen\'etica.
\end{itemize}

\section{Elitismo y Ciclo Evolutivo}

El proceso evolutivo se estructura en ciclos o generaciones.

\subsection*{Elitismo}
\begin{itemize}
    \item Se retiene el mejor individuo de la generaci\'on actual y se garantiza su inclusi\'on en la siguiente generaci\'on.
    \item Esto asegura que la calidad de la soluci\'on nunca empeore a lo largo de las generaciones.
\end{itemize}

\subsection*{Ciclo Evolutivo}
\begin{itemize}
    \item Cada generaci\'on incluye la selecci\'on, el cruzamiento, la mutaci\'on y la incorporaci\'on del individuo de \'elite.
    \item La evoluci\'on se repite durante un n\'umero predefinido de generaciones.
\end{itemize}

\subsection*{Registro y An\'alisis}
\begin{itemize}
    \item Se almacena el historial del fitness y de las mejores soluciones.
    \item Esto facilita el an\'alisis del comportamiento del algoritmo y la generaci\'on de visualizaciones.
\end{itemize}

\section{Penalización Exterior}

Los métodos de penalización son técnicas para manejar restricciones en problemas de optimización. La idea principal es transformar un problema de optimización con restricciones en uno sin restricciones, modificando la función objetivo para penalizar las soluciones que no cumplen con las restricciones. En el método de penalización exterior, esto se logra agregando términos a la función objetivo que aumentan su valor cuando se violan las restricciones. [cite: 45, 46, 47, 48, 49, 50, 51]

La función objetivo modificada, denotada como \( F_P(x) \), se define como:

\[
F_P(x) = f(x) + \lambda P(x)
\]

Donde:

\begin{itemize}
    \item \( f(x) \) es la función objetivo original del problema.
    \item \( P(x) \) es la función de penalización, que cuantifica la violación de las restricciones.
    \item \( \lambda_P \) es el parámetro de penalización, que controla la severidad de la penalización.
\end{itemize}

La función de penalización \( P(x) \) se calcula de la siguiente manera:

\[
P(x) = \sum_{i=1}^{n} \max(0, g_i(x))^2 + \sum_{j=1}^{m} h_j(x)^2
\]

Donde:

\begin{itemize}
    \item \( g_i(x) \) representa la \( i \)-ésima restricción de desigualdad, de la forma \( g_i(x) \leq 0 \).
    \item \( h_j(x) \) representa la \( j \)-ésima restricción de igualdad, de la forma \( h_j(x) = 0 \).
\end{itemize}

\subsection*{Funcionamiento de la Penalización Exterior}

\begin{itemize}
    \item Para las restricciones de desigualdad \( g_i(x) \leq 0 \):
    \begin{itemize}
        \item Si la restricción se cumple, \( \max(0, g_i(x)) = 0 \), por lo que no hay penalización.
        \item Si la restricción se viola, \( \max(0, g_i(x)) = g_i(x) \), y se aplica una penalización proporcional al cuadrado de la violación.
    \end{itemize}
    \item Para las restricciones de igualdad \( h_j(x) = 0 \):
    \begin{itemize}
        \item Si la restricción se cumple, \( h_j{(x)}^2 = 0 \), y no hay penalización.
        \item Si la restricción no se cumple, \( h_j{(x)}^2 > 0 \), y se aplica una penalización proporcional al cuadrado de la diferencia.
    \end{itemize}
\end{itemize}

\subsection*{Ventajas}

\begin{itemize}
    \item \textbf{Facilidad de implementación:} Es relativamente sencillo agregar la penalización externa a los algoritmos de optimización existentes.
    \item \textbf{Flexibilidad:} Permite explorar soluciones que violan las restricciones, lo cual puede ser útil en problemas complejos.
\end{itemize}

\subsection*{Desventajas}

\begin{itemize}
    \item \textbf{Distorsión de la solución óptima:} La penalización puede llevar a soluciones subóptimas si se priorizan soluciones que violan las restricciones pero tienen un mejor valor de la función objetivo. [cite: 61, 62]
    \item \textbf{Sensibilidad al parámetro de penalización:} La elección del valor de \( \lambda_P \) es crucial. Si es muy pequeño, las restricciones pueden no cumplirse; si es muy grande, puede dificultar la convergencia. [cite: 63, 64, 65]
    \item \textbf{Convergencia:} Puede causar una convergencia lenta o incluso la no convergencia del algoritmo.
\end{itemize}
\chapter{Resultados y Discusi\'on}

Durante la ejecución del algoritmo genético, se realizaron múltiples corridas completas para cada uno de los tres problemas de optimización de carteras, lo que permitió evaluar la estabilidad y eficiencia del método. Los resultados se agruparon en resúmenes globales, donde se registraron indicadores clave, los cuales se encuentran dentro de las siguientes tablas:

\foreach \i in {1,...,3}{
    \begin{longtable}{lrrrrrrr}
\caption{Problema 01: resumen\_global\_corridas}\label{tab:resumen_global_corridas} \\
\toprule
Indicador & x0 & x1 & x2 & x3 & x4 & x5 & Fitness \\
\midrule
\endfirsthead
\toprule
Indicador & x0 & x1 & x2 & x3 & x4 & x5 & Fitness \\
\midrule
\endhead
\midrule
\multicolumn{8}{r}{Continued on next page} \\
\midrule
\endfoot
\bottomrule
\endlastfoot
Mejor (Fitness) & 0.0 & 0.0565 & 0.4 & 0.3714 & 0.1746 & 0.0 & -0.6891 \\
Peor (Fitness) & 0.1285 & 0.0225 & 0.4 & 0.0905 & 0.1456 & 0.214 & -0.6114 \\
Media & 0.0484 & 0.0641 & 0.4 & 0.1991 & 0.2152 & 0.0749 & -0.6573 \\
Desv. Estándar & 0.0595 & 0.0349 & 0.0 & 0.1088 & 0.0952 & 0.0767 & 0.028 \\
\end{longtable}
}

Donde cada Indicador Representa lo siguiente:
\begin{itemize}
    \item \textbf{Mejor (Fitness):} Representa la soluci\'on con el valor de fitness m\'inimo obtenido en todas las corridas.
    \item \textbf{Peor (Fitness):} Indica la soluci\'on con el mayor valor de fitness, sirviendo como referencia de la variabilidad en la b\'usqueda.
    \item \textbf{Media:} Es el promedio de los valores de fitness de la mejor soluci\'on de cada corrida, ofreciendo una visi\'on global del desempe\~no del algoritmo.
    \item \textbf{Desv. Est\'andar:} Mide la dispersi\'on de los valores de fitness entre las corridas, reflejando la estabilidad y consistencia del proceso evolutivo.
\end{itemize}

\section{Análisis de los Resultados}

\subsection{Consistencia y Robustez}

Los resúmenes globales muestran que, a lo largo de las corridas, el algoritmo tiende a converger de manera consistente hacia soluciones de alta calidad. Una baja desviación estándar en los valores de fitness sugiere que el proceso evolutivo es robusto y no depende en exceso de la aleatoriedad inherente a los operadores genéticos. Esto es crucial para problemas de optimización de carteras, ya que garantiza que la metodología aplicada es reproducible y confiable.

\subsection{Análisis por Problema}

\subsubsection{Problema 1: Maximización de Retorno}

El resumen global para el Problema 1 indica la capacidad del algoritmo para identificar carteras con un alto retorno esperado. El valor de ``Mejor (Retorno)'' obtenido se sitúa en un rango competitivo, y la media de las corridas respalda la eficacia del algoritmo genético para explorar el espacio de soluciones y encontrar combinaciones de activos que maximizan el retorno, respetando la restricción de inversión máxima por activo.

\subsubsection{Problema 2: Minimización de Riesgo}

Para el Problema 2, los resultados globales evidencian una convergencia hacia carteras con un riesgo mínimo, cumpliendo con la restricción de rendimiento esperado mínimo. Los indicadores de ``Peor (Riesgo)'' y ``Desv. Estándar'' muestran que, aunque existen ciertas variaciones entre corridas, el algoritmo genético logra encontrar soluciones consistentes que equilibran el riesgo y el rendimiento esperado.

\subsubsection{Problema 3: Maximización de Retorno con Riesgo Controlado}

En el Problema 3, los resultados muestran la capacidad del algoritmo para maximizar el retorno esperado sin exceder un nivel de riesgo predefinido. Esto demuestra la efectividad del método de penalización para guiar la búsqueda hacia soluciones que satisfacen simultáneamente múltiples objetivos y restricciones.

\subsection{Evolución del Fitness y Visualizaciones}
Las gr\'aficas de la evoluci\'on del fitness, tanto en su forma original como normalizada, permiten observar el progreso generacional. Se aprecia una clara tendencia a la mejora, donde la mayor\'ia de las corridas muestran una reducci\'on significativa del valor de fitness a medida que avanzan las generaciones.

A continuación se dejan las gráficas de la evolución del fitness:
\foreach \i in {1,...,3}{
    \begin{figure}[H]
        \centering
        \includegraphics[width=\textwidth]{recursos/evolucion_fitness_with_constraints_Problem\i.png}
        \caption{Evolución del \text{fitness} del Problema \i}%
        \label{fig:fitness_problem_\i}
    \end{figure}
}
\section{Discusión Resultados}

\begin{itemize}
    \item \textbf{Eficacia del Algoritmo:} Los indicadores globales extra\'idos de los CSV demuestran que el algoritmo gen\'etico es capaz de acercarse a la soluci\'on \'optima, manteniendo una evoluci\'on progresiva y consistente en la reducci\'on del valor de fitness.

    \item \textbf{Diversidad y Convergencia:} La aplicaci\'on de operadores de selecci\'on, cruzamiento y mutaci\'on, junto con el mecanismo de elitismo, garantiza un equilibrio entre la exploraci\'on y la explotaci\'on del espacio de b\'usqueda. Esto se refleja en la baja variabilidad entre corridas, lo que es un indicativo de la estabilidad del proceso.

    \item \textbf{Potencial de Adaptaci\'on:} La estructura modular y la robustez mostrada por los resultados permiten considerar la posibilidad de aplicar este marco a otros problemas de optimizaci\'on, incluso aquellos con mayores dimensiones o con funciones objetivo de mayor complejidad.
\end{itemize}


\chapter{Implementaci\'on}

El proyecto se ha desarrollado siguiendo una arquitectura modular\footnote{Se puede encontrar el codigo perteneciente a cada funcion dentro de} que permite separar claramente las distintas funcionalidades y facilita tanto el mantenimiento como la ampliaci\'on futura. A continuaci\'on, se detallan los principales componentes y c\'omo se integran en el sistema:

\section{Funciones Objetivo}

\subsection{Descripci\'on}
Las funciones objetivo definen el problema a optimizar. En este proyecto se incluyen dos funciones:
\begin{itemize}
    \item \textbf{Tasa de retorno del portafolio}: Una funci\'on que nos permite calcular la ganancia esperada de un portafolio de inversion considerando las ganancias individuales de cada instrumento.
    \item \textbf{Varianza del portafolio}: Una funci\'on que nos permite calcular el riesgo de un portafolio de inversion considerando la covarianza de nuestros instrumentos de inversion.
\end{itemize}

\subsection{Implementaci\'on}
Se abordaron tres problemas de inversión, cada uno con sus respectivas limitaciones. En los tres casos se establecieron límites para todos los instrumentos de inversión en el intervalo \([0, 10]\). 

En el primer problema, se buscó maximizar el beneficio del portafolio, es decir, maximizar la tasa de retorno. En el segundo problema, el objetivo fue minimizar la varianza del portafolio, manteniendo una ganancia mínima del \(35\%\); por lo tanto, se incorporó esta restricción al modelo. Finalmente, en el tercer problema se volvió a buscar la maximización de la tasa de retorno del portafolio, pero añadiendo la restricción de considerar un riesgo menor al \(20\%\).

\section{M\'odulos del Algoritmo Gen\'etico}

El n\'ucleo del algoritmo gen\'etico se distribuye en varios m\'odulos:

\subsection{Inicializaci\'on}
\textbf{Funci\'on}: \texttt{initialize\_population}  
\textbf{Ubicaci\'on}: \texttt{libs/auxiliaries\_functions.py}  
\textbf{Descripci\'on}: Genera la poblaci\'on inicial de manera uniforme en el espacio de b\'usqueda.

\subsection{Selecci\'on}
\textbf{Funci\'on}: \texttt{vectorized\_tournament\_selection}  
\textbf{Ubicaci\'on}: \texttt{libs/selection.py}  
\textbf{Descripci\'on}: Se usa un enfoque de torneos para la selecci\'on de padres, empleando operaciones vectorizadas con NumPy.

\subsection{Cruzamiento}
\textbf{Funciones}: \texttt{sbx\_crossover}, \texttt{sbx\_crossover\_with\_boundaries}  
\textbf{Ubicaci\'on}: \texttt{libs/crossover.py}  
\textbf{Descripci\'on}: Implementa el operador SBX (Simulated Binary Crossover) con y sin control de l\'imites.

\subsection{Mutaci\'on}
\textbf{Funciones}: \texttt{polynomial\_mutation}, \texttt{polynomial\_mutation\_with\_boundaries}  
\textbf{Ubicaci\'on}: \texttt{libs/mutation.py}  
\textbf{Descripci\'on}: Aplica mutaci\'on polinomial, con control opcional de l\'imites para mantener la viabilidad de las soluciones.

\subsection{Aplicacion de Restricciones}
\textbf{Función}: \texttt{evaluate\_individuals\_with\_constraints} 
\textbf{Ubicaci\'on}: \texttt{libs/auxiliaries\_functions.py}  
\textbf{Descripción}: Verifica el cumplimiento de las restricciones impuestas por el problema. Cuando se violan restricciones, se aplica un mecanismo de penalización para ajustar el valor del fitness, de modo que las soluciones inviables tengan menor probabilidad de ser seleccionadas.

\subsection{Ejecuci\'on del Algoritmo}
\textbf{Funci\'on}: \texttt{genetic\_algorithm}  
\textbf{Ubicaci\'on}: \texttt{AG.py}  
\textbf{Descripci\'on}: Gestiona el ciclo evolutivo completo del algoritmo gen\'etico.

\section{Visualizaci\'on y Almacenamiento}

\subsection{Visualizaci\'on}
\textbf{M\'odulo}: \texttt{libs/plot.py}  
\textbf{Funciones}: \texttt{plot\_evolucion\_fitness}, \texttt{plot\_surface\_3d}  
\textbf{Descripci\'on}: Permite analizar la evoluci\'on del fitness y visualizar la superficie de las funciones objetivo.

\subsection{Almacenamiento}
\textbf{Estructura de Carpetas}:
\begin{itemize}
    \item Directorio \texttt{outputs} organizado en subcarpetas por problema.
    \item Historiales en archivos CSV con datos de fitness y variables.
    \item Res\'umenes estad\'isticos de cada corrida.
\end{itemize}

\textbf{Integraci\'on}: \texttt{main\_script.py} ejecuta el algoritmo para cada funci\'on definida en \texttt{AG\_confs.py}.

\textbf{Escalabilidad}: La arquitectura modular permite agregar nuevas funciones objetivo y modificar operadores gen\'eticos sin afectar la estructura base.
\chapter{Conclusiones}
Los resultados muestran que, bajo restricciones realistas como límites de inversión por activo y control del riesgo, es posible construir carteras 
bien diversificadas con alto potencial de rendimiento. La optimización permitió alinear expectativas de ganancia con distintos niveles de tolerancia 
al riesgo, haciendo uso eficiente de la información disponible.\\

El análisis confirma que considerar covarianzas y varianzas, además de los rendimientos esperados, es clave para evitar decisiones miopes. 
Así, se demuestra el valor práctico de los modelos de optimización en la construcción de portafolios robustos y alineados con objetivos claros.

\appendix
\chapter{GitHub}
Se anexa un enlace al repositorio donde se encuentran los codigos y una README que examnde la información aqui mostrada: \href{https://github.com/CARBAJE/Practica02-TSAB}{\textit{Click aqui para abrir el enlace al repositorio}}
funcion\chapter{Scripts}

\section{Archivo \texttt{main.py}}
\begin{lstlisting}[
  caption={Implementación de \texttt{main.py}},
  style=pythonstyle,
  basicstyle=\ttfamily\footnotesize
]
import os
import pandas as pd
from AG_confs import *

from AG import genetic_algorithm
from libs.plot import *

def main():
    # Crear carpetas de salida generales
    os.makedirs("outputs", exist_ok=True)
    
    for func_key, func_data in FUNCTIONS.items():
        f_obj = func_data["func"]
        lb = func_data["lb"]
        ub = func_data["ub"]
        func_name = func_data["name"]
        num_runs = func_data["num_runs"]
        
        # Carpetas especificas de cada funcion
        func_folder = f"outputs/{func_key}"
        os.makedirs(func_folder, exist_ok=True)
        hist_folder = os.path.join(func_folder, "historiales")
        res_folder = os.path.join(func_folder, "resumenes")
        os.makedirs(hist_folder, exist_ok=True)
        os.makedirs(res_folder, exist_ok=True)
        
        print(f"\n==============================================")
        print(f"  FUNCIoN: {func_name}")
        print(f"==============================================")
        
        all_runs_history = []
        best_solutions_all_runs = []  # Guardaremos los mejores individuos (x1, x2) de cada corrida
        best_values_across_runs = []  # Guardaremos el best_val (fitness) de cada corrida
        
        for run in range(num_runs):
            print(f"\nEjecucion {run+1}/{num_runs}")
            
            (best_sol, best_val,
             worst_sol, worst_val,
             avg_sol,  avg_val,
             std_val,
             best_fitness_history,
             best_x1_history,
             best_x2_history,
             population_final,
             fitness_final,
             best_solutions_over_time) = genetic_algorithm(
                 f_obj, lb, ub,
                 pop_size=POP_SIZE,
                 num_generations=NUM_GENERATIONS,
                 tournament_size=TOURNAMENT_SIZE,
                 crossover_prob=CROSSOVER_PROB,
                 eta_c=ETA_C,
                 mutation_prob=MUTATION_PROB,
                 eta_mut=ETA_MUT
             )
            
            # 1) Guardar historial
            df_historial = pd.DataFrame({
                "Generacion": np.arange(1, NUM_GENERATIONS + 1),
                "Mejor x1": best_x1_history,
                "Mejor x2": best_x2_history,
                "Mejor Fitness": best_fitness_history
            })
            historial_filename = os.path.join(hist_folder, f"historial_run_{run+1}.csv")
            df_historial.to_csv(historial_filename, index=False)
            
            # 2) Guardar resumen de la corrida
            data_resumen = [
                ["Mejor", best_sol[0], best_sol[1], best_val],
                ["Media", avg_sol[0], avg_sol[1], avg_val],
                ["Peor", worst_sol[0], worst_sol[1], worst_val],
                ["Desv. estandar", np.nan, np.nan, std_val]
            ]
            df_resumen = pd.DataFrame(data_resumen, columns=["Indicador", "x1", "x2", "Fitness"])
            resumen_filename = os.path.join(res_folder, f"resumen_run_{run+1}.csv")
            df_resumen.to_csv(resumen_filename, index=False)
            
            print(df_resumen.to_string(index=False))
            
            all_runs_history.append(best_fitness_history)
            best_solutions_all_runs.append(best_sol)
            best_values_across_runs.append(best_val)
        
        # ===========================================
        #       RESUMEN GLOBAL DE LAS CORRIDAS
        # ===========================================
        best_values_arr = np.array(best_values_across_runs)
        solutions_arr = np.array(best_solutions_all_runs)  # Cada fila: [x1, x2] del mejor individuo de cada corrida
        
        # Para el "Mejor" y "Peor", buscamos el indice de la corrida con minimo y maximo fitness
        min_index = np.argmin(best_values_arr)
        max_index = np.argmax(best_values_arr)
        
        data_global = [
            ["Mejor (Fitness)", solutions_arr[min_index, 0], solutions_arr[min_index, 1], best_values_arr[min_index]],
            ["Peor (Fitness)", solutions_arr[max_index, 0], solutions_arr[max_index, 1], best_values_arr[max_index]],
            ["Media", np.mean(solutions_arr[:, 0]), np.mean(solutions_arr[:, 1]), np.mean(best_values_arr)],
            ["Desv. Estandar", np.std(solutions_arr[:, 0]), np.std(solutions_arr[:, 1]), np.std(best_values_arr)]
        ]
        df_global = pd.DataFrame(data_global, columns=["Indicador", "x1", "x2", "Fitness"])
        
        global_filename = os.path.join(res_folder, "resumen_global_corridas.csv")
        df_global.to_csv(global_filename, index=False)
        
        # (Opcional) Graficar evolucion del fitness de todas las corridas
        plot_evolucion_fitness(all_runs_history, func_key, func_name)
        
        # (Opcional) Graficar superficie 3D si la funcion es de 2 variables
        if len(lb) == 2:
            plot_surface_3d(f_obj, lb, ub, best_solutions_all_runs, func_key, func_name)

if __name__ == "__main__":
    main()

\end{lstlisting}

\section{Archivo \texttt{AG\_confs.py}}
\begin{lstlisting}[
  caption={Implementación de \texttt{AG\_confs.py}},
  style=pythonstyle,
  basicstyle=\ttfamily\footnotesize
]
import numpy as np
from libs.functions import langermann, drop_wave
# ---------------------------
# Parametros del algoritmo
# ---------------------------
POP_SIZE = 100            # Numero de individuos en la poblacion
NUM_GENERATIONS = 200     # Numero de generaciones
NUM_RUNS = 10             # Numero de ejecuciones completas (ciclos)

# Parametros de la funcion de Langermann
a = np.array([3, 5, 2, 1, 7])
b = np.array([5, 2, 1, 4, 9])
c = np.array([1, 2, 5, 2, 3])

# Parametros del torneo
TOURNAMENT_SIZE = 3  # Numero de individuos participantes en cada torneo

# Parametros del cruzamiento SBX
CROSSOVER_PROB = 0.9  # Probabilidad de aplicar cruzamiento
ETA_C = 15            # indice de distribucion para SBX

# Parametros de la mutacion polinomial
MUTATION_PROB = 10.0 / 2  # Probabilidad de mutar cada gen
ETA_MUT = 20              # indice de distribucion para mutacion polinomial

best_solutions_list = [] 
all_runs_history = []  # Para graficar luego

FUNCTIONS = {
    "langermann": {
        "func": langermann,
        "lb": np.array([0, 0]),
        "ub": np.array([10, 10]),
        "name": "Langermann",
        "num_runs": NUM_RUNS
    },
    "drop_wave": {
        "func": drop_wave,
        "lb": np.array([-5.12, -5.12]),
        "ub": np.array([5.12, 5.12]),
        "name": "Drop-Wave",
        "num_runs": NUM_RUNS
    }
}
\end{lstlisting}

\section{Archivo \texttt{AG.py}}
\begin{lstlisting}[
  caption={Implementación de \texttt{AG.py}},
  style=pythonstyle,
  basicstyle=\ttfamily\footnotesize
]
from AG_confs import *
from libs.selection import vectorized_tournament_selection
from libs.crossover import sbx_crossover_with_boundaries
from libs.mutation import polynomial_mutation_with_boundaries
from libs.auxiliaries_functions import initialize_population


# ---------------------------
# Funcion principal del GA
# ---------------------------
def genetic_algorithm(objective_func, lower_bound, upper_bound,
                      pop_size=POP_SIZE, num_generations=NUM_GENERATIONS,
                      tournament_size=TOURNAMENT_SIZE,
                      crossover_prob=CROSSOVER_PROB, eta_c=ETA_C,
                      mutation_prob=MUTATION_PROB, eta_mut=ETA_MUT):
    """
    Ejecuta el GA para la funcion objetivo dada y retorna:
      - best_solution, best_value
      - worst_solution, worst_value
      - avg_solution, avg_value
      - std_value (fitness)
      - best_fitness_history, best_x1_history, best_x2_history
      - population (final), fitness (final)
      - best_solutions_over_time (para animaciones)
    """
    num_variables = len(lower_bound)
    
    # 1) Inicializar poblacion
    population = initialize_population(pop_size, num_variables, lower_bound, upper_bound)
    fitness = np.array([objective_func(ind) for ind in population])
    
    best_fitness_history = []
    best_x1_history = []
    best_x2_history = []
    
    # Para animacion: almacenamos el mejor (x1, x2) en cada generacion
    best_solutions_over_time = np.zeros((num_generations, num_variables))
    
    for gen in range(num_generations):
        # Elitismo: guardar el mejor de la generacion actual
        best_index = np.argmin(fitness)
        best_fitness = fitness[best_index]
        elite = population[best_index].copy()
        
        best_fitness_history.append(best_fitness)
        best_x1_history.append(elite[0])
        best_x2_history.append(elite[1])
        best_solutions_over_time[gen, :] = elite
        
        new_population = []
        
        # Numero de padres necesarios (2 por cada par a generar)
        num_parents_needed = 2 * (pop_size - 1)
        winners, _ = vectorized_tournament_selection(fitness, num_parents_needed,
                                                     tournament_size, len(population),
                                                     unique_in_column=True, unique_in_row=False)
        
        # Generar un valor global para el crossover y otro para la mutacion (para toda la generacion)
        global_u = np.random.rand()
        global_r = np.random.rand()
        
        # Generar nueva poblacion
        for i in range(0, len(winners), 2):
            parent1 = population[winners[i]].copy()
            if i + 1 < len(winners):
                parent2 = population[winners[i+1]].copy()
            else:
                parent2 = parent1.copy()
            
            # Cruzamiento SBX usando el mismo u para todas las variables del cruce
            child1, child2 = sbx_crossover_with_boundaries(
                parent1, parent2, lower_bound, upper_bound,
                eta_c, crossover_prob, use_global_u=True, global_u=global_u
            )
            # Mutacion polinomial usando el mismo r para todas las variables del individuo
            child1 = polynomial_mutation_with_boundaries(
                child1, lower_bound, upper_bound,
                mutation_prob, eta_mut, use_global_r=True, global_r=global_r
            )
            child2 = polynomial_mutation_with_boundaries(
                child2, lower_bound, upper_bound,
                mutation_prob, eta_mut, use_global_r=True, global_r=global_r
            )
            
            new_population.append(child1)
            if len(new_population) < pop_size - 1:
                new_population.append(child2)
        
        # Convertir a array y evaluar el fitness de la nueva poblacion
        new_population = np.array(new_population)
        new_fitness = np.array([objective_func(ind) for ind in new_population])
        
        # Incorporar el individuo elite (elitismo)
        new_population = np.vstack([new_population, elite])
        new_fitness = np.append(new_fitness, best_fitness)
        
        # Actualizar la poblacion y su fitness para la siguiente generacion
        population = new_population.copy()
        fitness = new_fitness.copy()
    
    # Calcular estadisticas finales
    best_index = np.argmin(fitness)
    worst_index = np.argmax(fitness)
    best_solution = population[best_index]
    best_value = fitness[best_index]
    worst_solution = population[worst_index]
    worst_value = fitness[worst_index]
    avg_solution = np.mean(population, axis=0)
    avg_value = np.mean(fitness)
    std_value = np.std(fitness)
    
    return (best_solution, best_value,
            worst_solution, worst_value,
            avg_solution, avg_value,
            std_value,
            best_fitness_history,
            best_x1_history,
            best_x2_history,
            population,
            fitness,
            best_solutions_over_time)
\end{lstlisting}

\section{Archivo \texttt{selection.py}}
\begin{lstlisting}[
  caption={Implementación de \texttt{selection.py}},
  style=pythonstyle,
  basicstyle=\ttfamily\footnotesize
]
import numpy as np

def vectorized_tournament_selection(fitness, num_tournaments, tournament_size, pop_size,
                                    unique_in_column=True, unique_in_row=False):
    """
    Genera una matriz de torneos de forma vectorizada y retorna, para cada torneo,
    el indice del individuo ganador (el de menor fitness).
    
    Args:
      - fitness: array con los fitness de la poblacion (longitud = pop_size).
      - num_tournaments: numero de torneos a realizar (por ejemplo, el numero total
                         de selecciones de padres requeridas en la generacion).
      - tournament_size: numero de individuos que participan en cada torneo.
      - pop_size: tamano de la poblacion.
      - unique_in_column: si True, para cada posicion (columna) se eligen candidatos sin
                          repeticion entre torneos.
      - unique_in_row: si True, en cada torneo (fila) los candidatos seran unicos.
                    (Por defecto se permite repetir en la fila).
    
    Returns:
      - winners: array de indices ganadores (uno por torneo).
      - tournament_matrix: la matriz de candidatos (de tamano [num_tournaments x tournament_size]).
    """
    if unique_in_row:
        # Para cada torneo (fila), muestreamos sin reemplazo (cada fila es unica)
        tournament_matrix = np.array([np.random.choice(pop_size, size=tournament_size, replace=False)
                                      for _ in range(num_tournaments)])
    else:
        # Permitir repeticion en la fila, pero controlar la no repeticion en cada columna
        if unique_in_column:
            # Para cada columna, se genera una permutacion de los indices (o se usan numeros aleatorios sin repeticion)
            # Siempre que num_tournaments <= pop_size.
            if num_tournaments > pop_size:
                # Si se requieren mas torneos que individuos, se hace sin la restriccion por columna.
                tournament_matrix = np.random.randint(0, pop_size, size=(num_tournaments, tournament_size))
            else:
                cols = []
                for j in range(tournament_size):
                    # Para la columna j, se toman num_tournaments indices sin repeticion
                    perm = np.random.permutation(pop_size)
                    cols.append(perm[:num_tournaments])
                tournament_matrix = np.column_stack(cols)
        else:
            # Sin restricciones, se muestrea con reemplazo para cada candidato.
            tournament_matrix = np.random.randint(0, pop_size, size=(num_tournaments, tournament_size))
    
    # Para cada torneo (fila de la matriz), se selecciona el candidato con el menor fitness.
    winners = []
    for row in tournament_matrix:
        row_fitness = fitness[row]
        winner_index = row[np.argmin(row_fitness)]
        winners.append(winner_index)
    winners = np.array(winners)
    return winners, tournament_matrix
\end{lstlisting}

\section{Archivo \texttt{crossover.py}}
\begin{lstlisting}[
  caption={Implementación de \texttt{crossover.py}},
  style=pythonstyle,
  basicstyle=\ttfamily\footnotesize
]
import numpy as np

def sbx_crossover(parent1, parent2, lower_bound, upper_bound, eta, crossover_prob):
    """Realiza el cruzamiento SBX para dos padres y devuelve dos hijos."""
    child1 = np.empty_like(parent1)
    child2 = np.empty_like(parent2)
    
    if np.random.rand() <= crossover_prob:
        for i in range(len(parent1)):
            u = np.random.rand()
            if u <= 0.5:
                beta = (2*u)**(1/(eta+1))
            else:
                beta = (1/(2*(1-u)))**(1/(eta+1))
            
            # Genera los dos hijos
            child1[i] = 0.5*((1+beta)*parent1[i] + (1-beta)*parent2[i])
            child2[i] = 0.5*((1-beta)*parent1[i] + (1+beta)*parent2[i])
            
            # Asegurar que los hijos esten dentro de los limites
            child1[i] = np.clip(child1[i], lower_bound[i], upper_bound[i])
            child2[i] = np.clip(child2[i], lower_bound[i], upper_bound[i])
    else:
        child1 = parent1.copy()
        child2 = parent2.copy()
    
    return child1, child2

def sbx_crossover_with_boundaries(parent1, parent2, lower_bound, upper_bound,
                                  eta, crossover_prob, use_global_u=False, global_u=None):
    """
    Realiza el cruzamiento SBX con limites, usando formulas que ajustan beta en funcion
    de la cercania a las fronteras. Permite usar un unico 'u' global para todos los individuos 
    de la generacion o, de forma estandar, un 'u' distinto por cada gen.
    
    Args:
      - parent1, parent2: arrays con los padres.
      - lower_bound, upper_bound: arrays con los limites inferiores y superiores.
      - eta: indice de distribucion para SBX.
      - crossover_prob: probabilidad de aplicar el cruce.
      - use_global_u: si es True se utilizara el mismo valor de 'u' para todas las variables.
      - global_u: valor de 'u' que se aplicara globalmente (si se proporciona).
      
    Returns:
      - child1, child2: arrays con los hijos resultantes.
    """
    parent1 = np.asarray(parent1)
    parent2 = np.asarray(parent2)
    child1 = np.empty_like(parent1)
    child2 = np.empty_like(parent2)
    
    # Si no se realiza el crossover, retornamos copias de los padres.
    if np.random.rand() > crossover_prob:
        return parent1.copy(), parent2.copy()
    
    # Si se quiere usar un 'u' global y no se ha pasado, se genera uno.
    if use_global_u:
        if global_u is None:
            global_u = np.random.rand()
    
    for i in range(len(parent1)):
        x1 = parent1[i]
        x2 = parent2[i]
        lb = lower_bound[i]
        ub = upper_bound[i]
        
        # Aseguramos que x1 sea menor o igual que x2
        if x1 > x2:
            x1, x2 = x2, x1
        
        dist = x2 - x1
        if dist < 1e-14:
            child1[i] = x1
            child2[i] = x2
            continue
        
        # Calcular la minima distancia a las fronteras
        min_val = min(x1 - lb, ub - x2)
        if min_val < 0:
            min_val = 0
        
        beta = 1.0 + (2.0 * min_val / dist)
        alpha = 2.0 - beta**(-(eta+1))
        
        # Si se usa u global, se usa el mismo valor para cada variable
        if use_global_u:
            u = global_u
        else:
            u = np.random.rand()
        
        if u <= (1.0 / alpha):
            betaq = (alpha * u)**(1.0/(eta+1))
        else:
            betaq = (1.0 / (2.0 - alpha*u))**(1.0/(eta+1))
        
        # Calcular los hijos
        c1 = 0.5 * ((x1 + x2) - betaq * (x2 - x1))
        c2 = 0.5 * ((x1 + x2) + betaq * (x2 - x1))
        
        # Ajustar a los limites
        child1[i] = np.clip(c1, lb, ub)
        child2[i] = np.clip(c2, lb, ub)
    
    return child1, child2
\end{lstlisting}

\section{Archivo \texttt{mutation.py}}
\begin{lstlisting}[
  caption={Implementación de \texttt{mutation.py}},
  style=pythonstyle,
  basicstyle=\ttfamily\footnotesize
]
import numpy as np

def polynomial_mutation(child, lower_bound, upper_bound, mutation_prob, eta_mut):
    """Aplica mutacion polinomial a un hijo."""
    mutant = child.copy()
    for i in range(len(child)):
        if np.random.rand() < mutation_prob:
            r = np.random.rand()
            diff = upper_bound[i] - lower_bound[i]
            if r < 0.5:
                delta = (2*r)**(1/(eta_mut+1)) - 1
            else:
                delta = 1 - (2*(1-r))**(1/(eta_mut+1))
            mutant[i] = child[i] + delta * diff
            mutant[i] = np.clip(mutant[i], lower_bound[i], upper_bound[i])
    return mutant


def polynomial_mutation_with_boundaries(child, lower_bound, upper_bound,
                                        mutation_prob, eta_mut,
                                        use_global_r=False, global_r=None):
    """
    Aplica mutacion polinomial (con limites) a un vector 'child'.
    Puede usar un unico 'r' global para todas las variables (si use_global_r=True)
    o generar un 'r' distinto para cada variable.
    
    Args:
      - child : array-like
          Cromosoma (vector de decision) a mutar.
      - lower_bound, upper_bound : array-like
          Limites inferiores y superiores para cada variable.
      - mutation_prob : float
          Probabilidad de mutacion (en [0,1]) para cada variable.
      - eta_mut : float
          indice de distribucion para la mutacion.
      - use_global_r : bool
          Si True, se utiliza un unico valor 'r' para todas las variables.
      - global_r : float, opcional
          Valor de 'r' global a usar; si no se proporciona, se genera uno.
    
    Retutrns:
      - mutant : np.ndarray
          Nuevo vector mutado (manteniendo la dimension de 'child').
    """
    mutant = np.array(child, copy=True, dtype=float)
    num_vars = len(child)
    
    # Si se desea usar un 'r' global y no se ha proporcionado, se genera uno una sola vez.
    if use_global_r:
        if global_r is None:
            global_r = np.random.rand()
    
    for i in range(num_vars):
        # Decidir si mutar esta variable
        if np.random.rand() < mutation_prob:
            x = mutant[i]
            xl = lower_bound[i]
            xu = upper_bound[i]
            
            # Evitar division por cero si los limites son casi iguales
            if abs(xu - xl) < 1e-14:
                continue

            # d = distancia normalizada al limite mas cercano
            d = min(xu - x, x - xl) / (xu - xl)
            
            # Elegir r: global o individual para cada variable
            if use_global_r:
                r = global_r
            else:
                r = np.random.rand()

            nm = eta_mut + 1.0

            # Calcular delta_q segun el valor de r
            if r < 0.5:
                bl = 2.0 * r + (1.0 - 2.0 * r) * ((1.0 - d) ** nm)
                delta_q = (bl ** (1.0 / nm)) - 1.0
            else:
                bl = 2.0 * (1.0 - r) + 2.0 * (r - 0.5) * ((1.0 - d) ** nm)
                delta_q = 1.0 - (bl ** (1.0 / nm))
            
            # Calcular la nueva posicion y asegurarse que este dentro de los limites
            y = x + delta_q * (xu - xl)
            mutant[i] = np.clip(y, xl, xu)
    
    return mutant
\end{lstlisting}

\section{Archivo \texttt{auxiliares\_functions.py}}
\begin{lstlisting}[
  caption={Implementación de \texttt{auxiliares\_functions.py}},
  style=pythonstyle,
  basicstyle=\ttfamily\footnotesize
]
import numpy as np
# ---------------------------
# Funciones auxiliares del GA
# ---------------------------
def initialize_population(pop_size, num_variables, lower_bound, upper_bound):
    """Inicializa la poblacion uniformemente en el espacio de busqueda."""
    return np.random.uniform(low=lower_bound, high=upper_bound, size=(pop_size, num_variables))
\end{lstlisting}
\chapter{Registro de indicadores completo.}

\section{Problema 01}

\subsection{Resumenes}
\foreach \i in {1,...,5}{
    \input{secciones/tablas/Problem1/resumenes/resumen_run_\i.tex}
}

\section{Problema 02}

\subsection{Resumenes}
\foreach \i in {1,...,5}{
    \input{secciones/tablas/Problem2/resumenes/resumen_run_\i.tex}
}

\section{Problema 03}

\subsection{Resumenes}
\foreach \i in {1,...,5}{
    \input{secciones/tablas/Problem3/resumenes/resumen_run_\i.tex}
}
\end{document}


\end{document}
